\documentclass[book.tex]{subfiles}

\begin{document}
\chapter{Gaussian Integrals}
\label{chap:gaussian_integrals}
\noindent\rule{11cm}{0.4pt} \\
\textbf{Prerequisites:} \autoref{chap:manifold} \\
\textbf{Difficulty Level:} *** \\
\noindent\rule{11cm}{0.4pt}
\vspace{5mm}


Let's start by considering the Gaussian integral

\[ G = \int_{-\infty}^\infty e^{-\frac{1}{2}x^2} dx \]

How do we evaluate this integral? The trick that's usually presented is to square the integral then do a change of basis to radial coordinates.

\begin{align}
    G^2 &= \left ( \int_{-\infty}^\infty e^{-\frac{1}{2}x^2} dx \right )^2 \\
    &= \int_{-\infty}^\infty \int_{-\infty}^\infty e^{-\frac{1}{2}(x^2 + y^2)} dx dy \\
\end{align}
Let's apply the change of variable 
\begin{align}
    x &= r \cos \theta \\
    y &= r \sin \theta \\
\end{align}
We can then compute the determinant of the Jacobian of this change of variables as 
\begin{align}
\det \mathcal{J} &= \det \begin{pmatrix}
    \frac{\partial x}{\partial r} & \frac{\partial x}{\partial \theta} \\
    \frac{\partial y}{\partial r} & \frac{\partial y}{\partial \theta} \\
    \end{pmatrix} \\
    &= \det \begin{pmatrix}
    \cos\theta & -r\sin \theta \\
    \sin \theta & r \cos \theta \\
\end{pmatrix} \\
&= r \\
\end{align}
This means in practice that we transform the volume element $dx dy$ into $r dr d\theta$. Let's plug this back into our integral to obtain
\begin{align}
    G^2 &= \int_{0}^{2\pi} \int_{0}^\infty e^{-\frac{1}{2}r^2} r dr d\theta\\
    &= 2 \pi \int_{0}^\infty e^{-\frac{1}{2}r^2} r dr \\
\end{align}
Let's apply the change of variable $u = \frac{1}{2}r^2$. Then $du = r dr$ to obtain
\begin{align}
    G^2 &= 2 \pi \int_0^\infty e^{-u} du \\
    &= 2 \pi
\end{align}

Thus it follows that $G = \sqrt{2 \pi}$. So far, this doesn't look too interesting. What does this integral calculation have to do with integral programs? Well one consideration here is that now that we have this closed form calculation for this simple Gaussian integral, anytime we see this same Gaussian integral in an integral program, the underlying compiler can just replace it with the constant $\sqrt{2 \pi}$. This isn't terribly interesting; how often are we likely to see just this constant?

Things start to look a little more interesting when we start to evaluate slightly more general Gaussian integrals.

Let's now consider the slightly more general integral

\[ \int_{-\infty}^{\infty} e^{-\frac{1}{2} ax^2} \]

We can take the same tack as above and compute

\begin{align}
    G^2 &= \left ( \int_{-\infty}^\infty e^{-\frac{1}{2}ax^2} dx \right )^2 \\
    &= \int_{-\infty}^\infty \int_{-\infty}^\infty e^{-\frac{1}{2}a(x^2 + y^2)} dx dy \\
\end{align}

Let's apply the same change of variables as above.

\begin{align}
    G^2 &= \int_{0}^{2\pi} \int_{0}^\infty e^{-\frac{1}{2}ar^2} r dr d\theta\\
    &= 2 \pi \int_{0}^\infty e^{-\frac{1}{2}ar^2} r dr \\
\end{align}

Let's  apply the change of variable $u = \frac{1}{2}ar^2$. Then $du = ar dr$ to obtain
\begin{align}
    G^2 &= 2 \pi \int_0^\infty \frac{e^{-u}}{a} du \\
    &= \frac{2 \pi}{a}
\end{align}
Thus it follows that $G = \sqrt{\frac{2\pi}{a}}$. Let's now evaluate a slightly more complex integral again.

\[ \int_{-\infty}^{\infty} e^{-\frac{1}{2}ax^2 + Jx} dx \]

Here $J$ is another constant. To solve this integral, we will work to complete the square on the quadratic $-\frac{1}{2}ax^2 + Jx$

\begin{align}
    -\frac{1}{2}ax^2 + J x &= -\frac{1}{2}a \left (x^2 -2\left (\frac{J}{a} \right ) x \right) \\
    &= -\frac{1}{2}a \left (x^2 -2\left (\frac{J}{a} \right ) x + \left ( \frac{J}{a} \right )^2 \right ) + \frac{J^2}{2a} \\
    &= -\frac{1}{2}a \left (x -\left (\frac{J}{a} \right ) \right )^2 + \frac{J^2}{2a} \\
\end{align}

Let's apply the change of variable $y = x + \frac{J}{a}$. Then we see

\begin{align}
    -\frac{1}{2} a x^2 + J x &= -\frac{1}{2} a y^2 + \frac{J^2}{2a}
\end{align}
Let's plug this back into the source integral
\begin{align}
    G &= \int_{-\infty}^{\infty} e^{-\frac{1}{2}ax^2 + Jx} \\
    &= \int_{-\infty}^{\infty} e^{-\frac{1}{2}ay^2 + J^2/2a} dy \\
    &= \sqrt{\frac{2 \pi}{a}} e^{J^2/2a}
\end{align}

\end{document}
